% Informe técnico — Proyecto Nota 3: Big Data con Python y PySpark (versión en español)
% Compilar desde esta carpeta:  pdflatex informe.tex  (dos veces, para el índice)
\documentclass[11pt,a4paper]{article}

\usepackage[utf8]{inputenc}
\usepackage[T1]{fontenc}
\usepackage[spanish,es-noshorthands,es-tabla]{babel}
\usepackage{lmodern}
\usepackage[margin=2.5cm]{geometry}
\usepackage{graphicx}
\usepackage{booktabs}
\usepackage{tabularx}
\usepackage{xcolor}
\usepackage{listings}
\usepackage{caption}
\usepackage{float}
\usepackage[hidelinks]{hyperref}

\graphicspath{{../../graficos/}}

\definecolor{codebg}{HTML}{F6F6F4}
\definecolor{codecomment}{HTML}{52514E}
\definecolor{pending}{HTML}{B3261E}

\lstset{
  language=Python,
  basicstyle=\ttfamily\small,
  backgroundcolor=\color{codebg},
  commentstyle=\color{codecomment}\itshape,
  keywordstyle=\bfseries,
  frame=single,
  rulecolor=\color{codebg},
  breaklines=true,
  columns=fullflexible,
  keepspaces=true,
  showstringspaces=false,
  captionpos=b,
  literate={á}{{\'a}}1 {é}{{\'e}}1 {í}{{\'i}}1 {ó}{{\'o}}1 {ú}{{\'u}}1 {ñ}{{\~n}}1
}
\renewcommand{\lstlistingname}{Fragmento}

% Marca visible para lo que falta antes de entregar.
\newcommand{\pendiente}[1]{\textcolor{pending}{\textbf{[PENDIENTE: #1]}}}

\title{Procesamiento de 66 millones de reseñas de Amazon:\\
pandas frente a Apache Spark (PySpark)\\[0.4em]
\large Proyecto Nota 3 — Big Data con Python y PySpark}
\author{Cristhian Eduardo Osorio Restrepo \and Daniel Stiven Pérez Córdoba \and Maria Camargo\\[0.3em]
\small Docente: Ing. Luis Fernando Castro Rojas, Mg, PhD.}
\date{Noviembre de 2026}

\begin{document}
\maketitle
\tableofcontents
\newpage

% ---------------------------------------------------------------------------
\section{Introducción y planteamiento del problema}

\subsection{Dominio y problema}

El comercio electrónico genera cada día millones de reseñas de productos. Una tienda o un
vendedor que quiera entender esa información necesita responder preguntas sencillas sobre
volúmenes muy grandes: cuándo se concentra la actividad, qué tiendas reciben más reseñas y
qué tan buenas son sus calificaciones, y qué registros están incompletos.

Este proyecto responde cinco preguntas de negocio sobre las reseñas de la categoría
\emph{Home and Kitchen} (hogar y cocina) de Amazon:

\begin{enumerate}
  \item ¿Cómo se comporta la cantidad de reseñas en el tiempo (año, mes y día de la semana)?
  \item ¿Qué tiendas concentran el mayor volumen de reseñas?
  \item ¿Cuál es la calificación media y su dispersión en cada grupo?
  \item ¿Qué registros son incompletos o incoherentes?
  \item ¿Qué reseñas tienen y cuáles no tienen su producto en el catálogo?
\end{enumerate}

Responderlas exige las operaciones que pide el enunciado: lectura, limpieza, columnas
derivadas, unión (\emph{join}) con el catálogo de productos, agregaciones por grupo y escritura
de resultados en Parquet.

\subsection{Datos}

Se usa el conjunto \emph{Amazon Reviews 2023} del laboratorio McAuley de la Universidad de
California en San Diego~\cite{hou2024}, publicado en Hugging Face. Se descarga por URL
directa, sin ejecutar código del repositorio de origen.

\begin{table}[H]
\centering
\caption{Archivos usados de la categoría \emph{Home and Kitchen}. Tamaños verificados con
las cabeceras HTTP el 17 de noviembre de 2026.}
\label{tab:datos}
\begin{tabularx}{\textwidth}{lXrr}
\toprule
Archivo & Contenido & Tamaño & Registros \\
\midrule
\texttt{0core/rating\_only} (CSV) & Interacciones usuario--producto: usuario, producto, calificación y fecha & 3,86 GB & 66\,623\,880 \\
\texttt{meta\_Home\_and\_Kitchen} (JSONL) & Catálogo de productos: categoría, tienda, precio & 11,8 GB & 3\,735\,584 \\
\texttt{review\_categories} (JSONL) & Reseñas completas, con título y texto & 31,4 GB & --- \\
\bottomrule
\end{tabularx}
\end{table}

Las mediciones de la Sección~\ref{sec:evidencia} se hicieron primero con el CSV de
interacciones, unido al catálogo. El archivo de reseñas completas, con el texto, se usa
para la prueba de volumen. \pendiente{confirmar con los resultados de la tarea T1.}

\subsection{¿Por qué es un problema de Big Data?}

\begin{itemize}
  \item \textbf{Volumen.} El archivo de reseñas completas pesa 31,4 GB, el doble de la
  memoria RAM del equipo de prueba (15,3 GB). Solo las interacciones suman 66,6 millones de
  filas.
  \item \textbf{Variedad.} Se combinan datos tabulares (CSV) con datos semiestructurados
  (JSON por línea) que incluyen texto libre.
  \item \textbf{Memoria, no disco.} Un archivo comprimido o columnar ocupa en disco mucho
  menos de lo que ocupa cargado en memoria. En las mediciones, unir las interacciones con el
  catálogo produce un \emph{DataFrame} de pandas de unos 12 GB, al límite de los 15,3 GB del
  equipo.
\end{itemize}

La hipótesis del proyecto es que, en un mismo equipo, pandas deja de ser viable cuando el
volumen supera la memoria disponible, mientras que Spark completa el procesamiento porque
trabaja por particiones.

% ---------------------------------------------------------------------------
\section{Marco teórico}

\subsection{Big Data}

Se habla de Big Data cuando el volumen, la velocidad de llegada o la variedad de los datos
superan lo que una herramienta tradicional procesa en un tiempo y con unos recursos
razonables. La respuesta habitual es \emph{escalar horizontalmente}: repartir los datos en
partes y procesarlas en paralelo, en varios núcleos o en varias máquinas, en lugar de
depender de una sola máquina más grande.

\subsection{Arquitectura de Apache Spark}

Apache Spark es un motor de procesamiento distribuido de código abierto~\cite{zaharia2016,chambers2018}.
Sus piezas principales son:

\begin{description}
  \item[Driver.] Es el proceso que ejecuta el programa del usuario. Construye el plan de
  ejecución, lo divide en tareas y las reparte. No debe recibir todos los datos: si se
  usa \texttt{collect()} sobre un conjunto grande, el driver se queda sin memoria.
  \item[Executors.] Son los procesos que ejecutan las tareas sobre los datos. En el modo
  local que usa este proyecto (\texttt{local[*]}), el driver y los executors viven en la
  misma máquina virtual de Java y hay un hilo de ejecución por núcleo.
  \item[Particiones.] Los datos se dividen en bloques llamados particiones. Cada tarea
  procesa una partición, de modo que nunca hace falta tener todo el conjunto en memoria a
  la vez. Si una partición no cabe, Spark escribe datos intermedios en disco (\emph{spill})
  en lugar de fallar.
  \item[Transformaciones y acciones.] Las transformaciones (\texttt{filter},
  \texttt{withColumn}, \texttt{join}, \texttt{groupBy}) describen un nuevo conjunto de datos
  sin calcularlo. Las acciones (\texttt{count}, \texttt{write}, \texttt{collect}) piden un
  resultado y disparan la ejecución.
  \item[Evaluación perezosa.] Como las transformaciones no se ejecutan de inmediato, Spark
  conoce el plan completo antes de correr. Eso le permite, por ejemplo, leer solo las
  columnas que se usan o filtrar antes de unir.
  \item[DAG.] El plan se representa como un grafo dirigido acíclico de operaciones. Spark lo
  corta en \emph{etapas} en cada punto donde los datos deben redistribuirse entre
  particiones (\emph{shuffle}), como un \texttt{groupBy} o un \texttt{join}. El shuffle es la
  operación más costosa, porque mueve datos entre tareas y los escribe en disco.
\end{description}

Sobre este núcleo, la API de DataFrames y Spark SQL~\cite{armbrust2015} agrega el optimizador
\emph{Catalyst}, que reescribe el plan lógico antes de ejecutarlo, y desde Spark~3 la
\emph{ejecución adaptativa} (AQE), que ajusta el plan con estadísticas reales, por ejemplo
uniendo particiones pequeñas~\cite{sparkdocs}.

\subsection{Técnicas de optimización usadas}

\begin{description}
  \item[Parquet.] Formato columnar y comprimido~\cite{parquet}. Al guardar cada columna por
  separado, una consulta lee solo las columnas que necesita (\emph{column pruning}).
  \item[Broadcast join.] Si una de las dos tablas es pequeña, se copia completa a cada
  executor y la tabla grande se une donde ya está, sin shuffle.
  \item[Caching.] Guarda en memoria un resultado intermedio que se reutiliza varias veces,
  para no recalcularlo. Solo conviene si recalcular cuesta más que guardar.
  \item[Particiones de shuffle.] Por defecto Spark crea 200 particiones en cada shuffle; en
  una sola máquina eso genera muchas tareas diminutas. Ajustarlas al número de núcleos
  reduce el costo de planificación.
\end{description}

\subsection{Limitaciones de pandas}

pandas~\cite{mckinney2010} es la biblioteca de referencia para análisis de datos tabulares
en Python, pero fue diseñada para un solo nodo y para datos que caben en memoria:

\begin{itemize}
  \item \textbf{Todo en memoria.} Carga el conjunto completo en RAM antes de operar. Si no
  cabe, el proceso termina con \texttt{MemoryError} o el sistema operativo lo cierra.
  \item \textbf{Ejecución inmediata.} Cada operación se ejecuta en el momento y crea
  resultados intermedios completos; un \texttt{merge} puede necesitar a la vez la tabla
  original y la unida.
  \item \textbf{Un solo núcleo.} La mayoría de operaciones usa un núcleo, aunque el equipo
  tenga varios.
\end{itemize}

Existen mitigaciones dentro de pandas: leer solo algunas columnas, usar tipos de datos más
pequeños o procesar el archivo por partes (\emph{chunks}). La propia documentación de pandas
las recomienda~\cite{pandasscale}. La Sección~\ref{sec:evidencia} muestra hasta dónde llegan.

% ---------------------------------------------------------------------------
\section{Entorno y configuración en Windows}

\subsection{Equipos}

Todas las mediciones se hicieron en un solo equipo Windows, con la misma configuración para
pandas y para Spark. El equipo Linux se usó solo para desarrollar y probar el código.

\begin{table}[H]
\centering
\caption{Equipos y versiones de software.}
\label{tab:entorno}
\begin{tabularx}{\textwidth}{lXX}
\toprule
 & Equipo de medición & Equipo de desarrollo \\
\midrule
Sistema operativo & Windows 11 Home & Debian 13 (Linux) \\
Procesador & AMD Ryzen 5 7520U, 4 núcleos y 8 hilos & 16 hilos \\
Memoria RAM & 15,3 GB & 15,3 GB \\
Disco libre & 300 GB & 13 GB \\
Python & 3.13.17 & 3.13.5 \\
Java & Microsoft OpenJDK 17.0.21 & Temurin 17.0.18 \\
\midrule
PySpark & \multicolumn{2}{l}{4.3.0 (incluye Spark)} \\
pandas / pyarrow / psutil & \multicolumn{2}{l}{3.1.1 / 26.0.1 / 7.2.2} \\
\bottomrule
\end{tabularx}
\end{table}

\subsection{Instalación}

La instalación completa está en el \texttt{README.md} del repositorio. En resumen:

\begin{enumerate}
  \item Python 3.13 con \texttt{winget}, usando siempre el lanzador \texttt{py -3.13}.
  \item JDK 17 con \texttt{winget}; el instalador define \texttt{JAVA\_HOME}.
  \item \texttt{winutils.exe} y \texttt{hadoop.dll} de Hadoop 3.3.6 con el script
  \texttt{scripts/setup\_winutils.ps1}, que verifica su SHA-256 y define
  \texttt{HADOOP\_HOME}.
  \item Un entorno virtual con \texttt{requirements.txt}, que fija las versiones exactas.
\end{enumerate}

\texttt{SPARK\_HOME} no es necesario: el paquete \texttt{pyspark} de PyPI trae Spark
completo y lo localiza por sí mismo. Se documenta así porque el enunciado menciona esa
variable.

\subsection{Dificultades encontradas y cómo se resolvieron}

\begin{table}[H]
\centering
\small
\caption{Problemas de entorno registrados en la bitácora del proyecto.}
\label{tab:dificultades}
\begin{tabularx}{\textwidth}{XX}
\toprule
Problema & Solución \\
\midrule
Solo estaba instalado Python 3.14, fuera del rango que soporta PySpark 4.3.0 (3.11--3.13), y el comando \texttt{python} abría la Microsoft Store. &
Instalar Python 3.13 con \texttt{winget} y usar siempre \texttt{py -3.13}. \\
No había Java ni \texttt{JAVA\_HOME}. &
Instalar Microsoft OpenJDK 17 con \texttt{winget}, que pide elevación y define la variable. \\
Escribir Parquet fallaba: primero por falta de \texttt{HADOOP\_HOME} y luego con \texttt{UnsatisfiedLinkError} por falta de \texttt{hadoop.dll}. &
Script que descarga \texttt{winutils.exe} y \texttt{hadoop.dll}, verifica su SHA-256 y define \texttt{HADOOP\_HOME} y \texttt{PATH} sin permisos de administrador. \\
El propio script no arrancaba en PowerShell 5.1: error de sintaxis en \texttt{"\$name: ..."} y acentos dañados al leer UTF-8 sin BOM. &
Usar \texttt{"\$\{name\}: ..."}, guardar con BOM y agregar una prueba automática que revisa la sintaxis del script. \\
PySpark 4.3.0 avisa que no soporta del todo pandas 3. &
Aviso sin efecto en las pruebas; se mantiene pandas 3.1.1 para no invalidar las mediciones ya hechas. \\
La descarga de archivos grandes no se podía reanudar: un corte de red obligaba a empezar de cero. &
\pendiente{describir la solución aplicada en la tarea T1.} \\
\bottomrule
\end{tabularx}
\end{table}

% ---------------------------------------------------------------------------
\section{Análisis del código fuente}

\subsection{Organización}

El repositorio separa lo que se entrega de lo que se mide:

\begin{itemize}
  \item \texttt{pandas\_version/analysis.py} y \texttt{spark\_version/analysis.py} contienen
  las mismas seis operaciones, una función por operación y con los mismos nombres.
  \item \texttt{spark\_version/optimized.py} reutiliza esas funciones de Spark y solo
  cambia \emph{cómo} se ejecutan.
  \item \texttt{benchmark/} importa esas mismas funciones para medirlas. Así, lo que se mide
  es exactamente lo que se entrega.
  \item \texttt{tests/} comprueba que pandas y Spark producen los mismos resultados sobre los
  mismos datos.
\end{itemize}

Las reglas comunes están escritas en un contrato (\texttt{docs/CONTRATO\_ANALISIS\_amazon\_hk.md}):
mismas columnas, mismas reglas de limpieza, misma numeración de días y tolerancia de
$10^{-4}$ en medias y desviaciones.

\subsection{Lectura}

Spark lee con un \textbf{esquema explícito}. Sin él, Spark tendría que recorrer el archivo
una vez solo para deducir los tipos de cada columna, lo que con decenas de gigabytes es una
pasada completa extra.

\begin{lstlisting}[caption={Esquema explícito del catálogo en Spark.}]
META_SCHEMA = StructType([
    StructField("parent_asin", StringType(), True),
    StructField("store", StringType(), True),
    StructField("price", DoubleType(), True),
    # ... otras columnas del catálogo
])
spark.read.schema(META_SCHEMA).json(str(meta_path)) \
    .filter(F.col("parent_asin").isNotNull()) \
    .dropDuplicates(["parent_asin"])
\end{lstlisting}

En la lectura no se carga nada todavía: es una transformación. pandas, en cambio, carga el
archivo completo en memoria en cuanto se llama a \texttt{read\_parquet}.

\subsection{Limpieza}

Las dos versiones aplican las mismas tres reglas: descartar filas con nulos en las columnas
clave, conservar calificaciones entre 1 y 5 y fechas positivas.

\begin{lstlisting}[caption={Limpieza en pandas (arriba) y en Spark (abajo).}]
cleaned = reviews.dropna(subset=["user_id", "parent_asin", "rating", "timestamp"])
cleaned = cleaned[(cleaned["rating"] >= 1) & (cleaned["rating"] <= 5)
                  & (cleaned["timestamp"] > 0)]

reviews.dropna(subset=["user_id", "parent_asin", "rating", "timestamp"]) \
    .filter((F.col("rating") >= 1) & (F.col("rating") <= 5) & (F.col("timestamp") > 0))
\end{lstlisting}

La sintaxis es casi igual, pero el efecto no: pandas crea en ese momento una copia filtrada
de toda la tabla, mientras que Spark solo agrega un paso al plan. Ese paso se ejecuta
después, partición por partición, y el optimizador lo ubica lo más cerca posible de la
lectura.

\subsection{Columnas derivadas}

Ambas versiones convierten la fecha (milisegundos desde 1970, en UTC) en año, mes, día de la
semana y hora. El detalle delicado es el \textbf{día de la semana}: pandas numera de 0
(lunes) a 6 (domingo) y Spark de 1 (domingo) a 7 (sábado). pandas se ajusta a Spark:

\begin{lstlisting}[caption={Conversión del día de la semana de pandas a la numeración de Spark.}]
derived["review_dow"] = ((review_ts.dt.dayofweek + 1) % 7) + 1
# lunes:   ((0 + 1) % 7) + 1 = 2   -> igual que Spark
# domingo: ((6 + 1) % 7) + 1 = 1   -> igual que Spark
\end{lstlisting}

Sin esta conversión, la tabla temporal saldría distinta en cada motor aunque ambos
estuvieran bien programados. Una prueba automática lo verifica para los siete días.

\subsection{Unión con el catálogo}

Es una unión \textbf{izquierda} por \texttt{parent\_asin}: toda reseña se conserva, y las que
no tienen producto en el catálogo quedan con la tienda vacía. No se descartan, porque
responden la pregunta~5.

\begin{lstlisting}[caption={Unión en pandas, en Spark base y en Spark optimizado.}]
reviews.merge(catalog, on="parent_asin", how="left")          # pandas
reviews.join(catalog, "parent_asin", "left")                  # Spark base
reviews.join(F.broadcast(catalog), "parent_asin", "left")     # Spark optimizado
\end{lstlisting}

En pandas, \texttt{merge} construye la tabla unida completa en memoria; es la operación que
más memoria necesita. En Spark base, la unión exige un shuffle de las dos tablas por
\texttt{parent\_asin}. En la versión optimizada el catálogo se reduce a dos columnas y se
difunde a cada executor (\emph{broadcast}), de modo que las reseñas se unen donde están, sin
moverlas.

\subsection{Agregaciones}

\begin{lstlisting}[caption={Agregado por tienda en pandas (arriba) y en Spark (abajo).}]
reviews.groupby("store", dropna=False).agg(
    reviews=("parent_asin", "size"),
    products=("parent_asin", "nunique"),
    avg_rating=("rating", "mean"),
    stddev_rating=("rating", "std"))

reviews.groupBy("store").agg(
    F.count("*").alias("reviews"),
    F.countDistinct("parent_asin").alias("products"),
    F.round(F.avg("rating"), 4).alias("avg_rating"),
    F.round(F.stddev("rating"), 4).alias("stddev_rating"))
\end{lstlisting}

Dos decisiones aseguran que den lo mismo:

\begin{itemize}
  \item \texttt{dropna=False} en pandas conserva el grupo de reseñas sin tienda. Spark
  conserva el grupo nulo por defecto; pandas lo descartaría sin esa opción.
  \item La desviación es \textbf{muestral} en ambos: \texttt{std} de pandas y
  \texttt{stddev} de Spark dividen entre $n-1$.
\end{itemize}

En Spark, un \texttt{groupBy} se ejecuta en dos fases: cada partición agrega sus propias
filas, se hace un shuffle de esos resultados parciales y luego se combinan. Por eso solo se
mueven los resultados parciales y no las filas originales.

\subsection{Escritura}

pandas escribe un archivo Parquet por tabla. Spark escribe una \emph{carpeta} con un archivo
por partición, porque cada tarea escribe lo suyo en paralelo. Como las tablas agregadas son
pequeñas, la versión optimizada aplica \texttt{coalesce(1)} para dejar un solo archivo por
tabla.

La escritura es además la \textbf{acción} que dispara todo el cálculo anterior. Por eso, en
las mediciones por operación, el tiempo de escritura de Spark incluye la lectura, la
limpieza, la unión y las agregaciones que hasta ese momento solo estaban planeadas.

\subsection{Versión optimizada y su relación con el DAG}

\texttt{spark\_version/optimized.py} reutiliza las funciones anteriores y cambia la ejecución
en cinco puntos, más un cache opcional. Cada uno se puede activar o apagar desde la línea de comandos para medir su efecto.

\begin{table}[H]
\centering
\small
\caption{Optimizaciones de la versión Spark y concepto que aplican.}
\label{tab:optimizaciones}
\begin{tabularx}{\textwidth}{lXl}
\toprule
Optimización & Qué cambia & Concepto \\
\midrule
Lectura de 4 columnas & Parquet no lee las demás columnas del disco & Formato columnar \\
Tabla temporal sin join & La tabla temporal no usa columnas del catálogo; se calcula sin la unión & DAG: un shuffle menos \\
Broadcast join & El catálogo, reducido a 2 columnas, se copia a cada executor & Broadcast join \\
Particiones de shuffle & 2 por núcleo en lugar de 200 & Particiones \\
Un archivo por tabla & \texttt{coalesce(1)} antes de escribir & Particiones \\
Cache (apagado) & Guardaría las reseñas limpias para las dos agregaciones & Caching \\
\bottomrule
\end{tabularx}
\end{table}

El efecto se ve en el plan físico que Spark genera antes de ejecutar (\texttt{explain}). Para
la tabla por tienda, la unión aparece como \texttt{BroadcastHashJoin} y la lectura solo pide
las columnas necesarias:

\begin{lstlisting}[language={},caption={Plan físico de la tabla por tienda (versión optimizada).}]
AdaptiveSparkPlan
+- HashAggregate                          <- combina los parciales
   +- Exchange                            <- shuffle (etapa nueva)
      +- HashAggregate                    <- agrega dentro de cada partición
         +- ...
            +- BroadcastHashJoin LeftOuter BuildRight
               :- Filter <- Scan parquet  (user_id, parent_asin, rating, timestamp)
               +- BroadcastExchange
                  +- Filter <- Scan parquet  (parent_asin, store)
\end{lstlisting}

El plan de la tabla temporal no contiene ningún \emph{join}: solo lectura, filtro,
proyección y agregación.

En una prueba con datos sintéticos de 20 millones de reseñas en el equipo Linux, la versión
optimizada produjo tablas idénticas a la base y tardó 25,7 a 26,1 segundos, frente a 27,1 a
27,7 de la base. Con el cache activado tardó más (28,8 a 29,2 segundos): releer cuatro
columnas de Parquet costó menos que guardar 20 millones de filas en memoria. Por eso el cache
queda apagado por defecto. \pendiente{reemplazar por la medición en el equipo Windows (tarea T4).}

% ---------------------------------------------------------------------------
\section{Evidencia de inviabilidad de pandas}
\label{sec:evidencia}

\pendiente{depende de la medición con el archivo de reseñas completas (tareas T1, T2 y T3):
errores de memoria, tiempos, capturas del Administrador de tareas y resultado de las
mitigaciones de pandas.}

% ---------------------------------------------------------------------------
\section{Análisis comparativo pandas frente a Spark}

\pendiente{tablas y gráficos con la medición final. Los gráficos actuales, con el CSV de
interacciones, están en \texttt{docs/graficos/}.}

% ---------------------------------------------------------------------------
\section{Conclusiones}

\pendiente{se escriben con los resultados de las secciones 5 y 6.}

% ---------------------------------------------------------------------------
\begin{thebibliography}{9}

\bibitem{hou2024}
Y.~Hou, J.~Li, X.~Fu, Z.~He, A.~Yan, X.~Chen y J.~McAuley, ``Bridging language and items
for retrieval and recommendation: Benchmarking LLMs as semantic encoders,'' \emph{arXiv
preprint arXiv:2403.03952v2}, 2024.

\bibitem{zaharia2016}
M.~Zaharia \emph{et al.}, ``Apache Spark: A unified engine for big data processing,''
\emph{Communications of the ACM}, vol.~59, no.~11, pp.~56--65, 2016.

\bibitem{armbrust2015}
M.~Armbrust \emph{et al.}, ``Spark SQL: Relational data processing in Spark,'' en
\emph{Proc. ACM \mbox{SIGMOD} International Conference on Management of Data}, 2015, pp.~1383--1394.

\bibitem{mckinney2010}
W.~McKinney, ``Data structures for statistical computing in Python,'' en \emph{Proc. 9th
Python in Science Conference (SciPy)}, 2010, pp.~56--61.

\bibitem{chambers2018}
B.~Chambers y M.~Zaharia, \emph{Spark: The Definitive Guide: Big Data Processing Made
Simple}. Sebastopol, CA, EE.~UU.:
O'Reilly Media, 2018.

\bibitem{sparkdocs}
The Apache Software Foundation, ``Spark SQL and DataFrames,'' documentación de Apache
Spark 4.3.0. [En línea]. Disponible:
\url{https://spark.apache.org/docs/4.3.0/sql-programming-guide.html}. Accedido: 17 de
noviembre de 2026.

\bibitem{parquet}
The Apache Software Foundation, ``Documentation,'' Apache Parquet. [En línea]. Disponible:
\url{https://parquet.apache.org/docs/}. Accedido: 17 de noviembre de 2026.

\bibitem{pandasscale}
pandas development team, ``Scaling to large datasets,'' documentación de pandas 3.1.1.
[En línea]. Disponible: \url{https://pandas.pydata.org/docs/user_guide/scale.html}.
Accedido: 17 de noviembre de 2026.

\end{thebibliography}

\end{document}
